\chapter{Online Assessments}\label{ch:iv:online-assessments}

A timer in the corner of the screen counts down from eight minutes; the counter reads question 1 of 80; a
wrong answer costs a mark and a blank costs nothing. The candidate who answers everything finishes with a lower
expected score than the one who skipped the fifth of the questions she was least sure of, and the test was
designed so that she would. An \omterm{def:iv:online-assessments:oa}{online assessment} measures a small number of abilities under a scoring rule,
and the rule is part of the question.

\section{Timed arithmetic and the scoring rule}

\begin{definition}[Online assessment, numerical reasoning test]\label{def:iv:online-assessments:oa}
An \emph{online assessment}\index{online assessment} is a standardised test taken remotely, under a time
limit, early in a process: arithmetic, numerical or logical reasoning, probability, coding, or a game. A
\emph{numerical reasoning test}\index{numerical reasoning test} is an online assessment of calculation with
numbers, fractions and percentages, or of reading quantities from tables and charts, scored for the number of
correct answers in the time.
\end{definition}

Two things decide the score besides ability: speed, and what the rule does with wrong answers. Eighty items in
eight minutes is six seconds an item, which is enough for a well-practised method
(\cref{ch:iv:mental-arithmetic}) and not enough for long multiplication. When wrong answers are penalised,
the rule also decides which items to answer.

\begin{proposition}[When to answer]\label{prop:iv:online-assessments:answer}
If a correct answer scores $+1$, a wrong one $-w$ and a blank 0, answering an item that you get right with
probability $p$ is worth $p - w(1-p)$ in expectation, which is positive exactly when $p > w/(1+w)$.
\end{proposition}

With $w = 1$ the threshold is one half; with the classical quarter-mark penalty on five-option items it is
$\tfrac15$, which is exactly the chance of a random guess, so that blind guessing is worth nothing and any
elimination of an option makes guessing worthwhile (\cref{fig:iv:online-assessments:skip}). The hook's
arithmetic: with 64 items answered at 95\% and 16 at 35\%, answering everything under $w = 1$ is worth $64
\times 0.90 + 16 \times (-0.30) = 52.8$, and skipping the sixteen weak items is worth 57.6.

\begin{omfigure}
\begin{tikzpicture}
\begin{axis}[width=10cm,height=5.2cm,xlabel={\small probability of a correct answer $p$},
  ylabel={\small expected marks per item},xmin=0,xmax=1,ymin=-1,ymax=1,
  tick label style={font=\footnotesize},grid=major,grid style={gray!20},table/col sep=comma,
  legend cell align=left,legend style={font=\scriptsize,at={(0.5,-0.32)},anchor=north,
  /tikz/every even column/.append style={column sep=8pt}},legend columns=3]
\addplot[omDef,thick] table[x=p,y=w0]{figdata/interviews/04-online-assessments/skip.csv};
\addplot[omProp,thick,dashed] table[x=p,y=w025]{figdata/interviews/04-online-assessments/skip.csv};
\addplot[omThm,thick,densely dotted] table[x=p,y=w1]{figdata/interviews/04-online-assessments/skip.csv};
\draw[gray] (axis cs:0,0) -- (axis cs:1,0);
\legend{no penalty,{penalty $\tfrac14$},penalty 1}
\end{axis}
\end{tikzpicture}

\omcaption{The value of answering an item against the chance of getting it right, for three penalties on a
wrong answer. Each line crosses zero at $p = w/(1+w)$: 0, 0.2 and 0.5. Data:
\texttt{fig\_iv\_skip.py}.}\label{fig:iv:online-assessments:skip}
\end{omfigure}

\begin{method}[Sitting a timed numerical test]\label{met:iv:online-assessments:timed}
\begin{enumerate}
\item Before starting, read the rules for three things: the penalty for a wrong answer, whether items can be
revisited, and whether difficulty rises through the test.
\item Set a time per item (the limit divided by the count) and a hard ceiling at about twice that; past the
ceiling, answer if above the threshold of \cref{prop:iv:online-assessments:answer} and move on.
\item If items can be skipped and revisited, take a first pass at every item you can do inside the time per
item, then return.
\item Practise with the same format, timed, until the methods of \cref{ch:iv:mental-arithmetic} are automatic.
Practice raises scores: a meta-analysis of 107 samples found an average gain of a quarter of a standard
deviation on retesting (Hausknecht et al., 2007), which is why many firms limit retakes
(\cref{dat:iv:the-process-by-firm-type:published}).
\end{enumerate}
\end{method}

\section{Sequences, patterns and logic items}

Sequence and logic items test whether the candidate can find a rule quickly and check it. Most number
sequences in such tests are built from a small set of rules, and trying them in a fixed order is faster than
staring.

\begin{method}[Finding the rule of a sequence]\label{met:iv:online-assessments:sequence}
\begin{enumerate}
\item Take differences; if they are constant after one or two rounds, the sequence is a polynomial of that
degree and the next term follows by adding back.
\item Take ratios; a constant ratio, or a ratio plus a constant ($a_{n+1} = 2a_n - 1$), is the next family.
\item Split the odd and even positions: two interleaved sequences are common.
\item Look for squares, cubes, primes, factorials and sums of the two previous terms.
\item Check the rule on every given term before answering.
\end{enumerate}
\end{method}

\begin{example}[Differences]\label{ex:iv:online-assessments:diff}
For 2, 6, 12, 20, 30 the differences are 4, 6, 8, 10 and the second differences are constant at 2: the next
difference is 12 and the next term 42. (The terms are $n(n+1)$, which is the check.)
\end{example}

Logic items are constraint problems: a few facts about an arrangement, one question about it. The method is
elimination: write the positions, place the most constrained item first, and test each remaining case against
every constraint. A candidate who writes the five seats on paper answers in a minute a question that takes
five in the head (\cref{iq:iv:online-assessments:7}).

\section{Automated coding assessments}

\begin{definition}[Automated coding assessment]\label{def:iv:online-assessments:coding}
An \emph{automated coding assessment}\index{automated coding assessment} is an \omterm{def:iv:online-assessments:oa}{online assessment} in which the
candidate writes programs that are run against hidden test cases, with limits on time and memory, and are
scored by the share of cases passed.
\end{definition}

Hidden tests have a predictable structure: small cases, edge cases (empty input, a single element, all
elements equal, the largest and smallest values), and large cases that only an efficient solution finishes in
the time limit. The input bounds in the statement say which complexity is expected.

\begin{method}[Reading the bounds]\label{met:iv:online-assessments:bounds}
\begin{enumerate}
\item At about $10^8$ simple operations a second: $n \le 10^3$ allows $O(n^2)$; $n \le 10^5$ to $10^6$ asks for
$O(n \log n)$ or $O(n)$; $n \le 20$ allows exponential search.
\item Write the edge cases down before coding, then code, then run them.
\item Watch integer overflow when products or sums of large values are asked for, and floating-point
equality when prices are decimals.
\item Submit a correct simple version early if partial credit is given, then improve.
\end{enumerate}
\end{method}

\section{Probability quizzes, games and the rules of the test}

Some assessments are short probability quizzes under a clock, testing the speed of the methods of
\cref{ch:iv:probability-i}; others are games that record behaviour (how much risk the candidate takes as a
reward grows, how quickly they learn a hidden rule) and score it against what the role needs. A game rarely
says what it scores; the useful assumption is that it scores what a good decision-maker would do, which can
usually be computed (\cref{iq:iv:online-assessments:8}).

Tests must be fair to candidates with disabilities, and candidates may ask for adjustments (extra time, a
different format) without that request being held against them.

\begin{dated}{2026-09}{Adjustments in recruitment tests}\label{dat:iv:online-assessments:adjust}
\textbf{United Kingdom}: section 20 of the Equality Act 2010 imposes a duty to make reasonable adjustments
where a provision, criterion or practice puts a disabled person at a substantial disadvantage; Schedule 8
applies it to employers, including their arrangements for deciding to whom to offer employment.
\textbf{United States}: under the regulations implementing the Americans with Disabilities Act (29 CFR
1630.11), an employer must select and administer employment tests so that, for an applicant with a disability
that impairs sensory, manual or speaking skills, the results reflect the ability the test purports to
measure rather than the impairment.
\end{dated}

Integrity rules matter as much: assessments forbid help from other people and, increasingly, from AI
assistants; a result obtained that way is found out in the next stage, which asks the same kind of question
live.

\section{Worked answers}

\begin{example}[Behind the clock at item 20]\label{ex:iv:online-assessments:behind}
\emph{Sixty items in thirty minutes; after item 20 the clock shows twelve minutes gone.} The budget was 30 seconds an
item, so item 20 should have been done at ten minutes: you are two minutes behind. The remaining 40 items now have
18 minutes, 27 seconds each. The plan that loses least is not to hurry every item by 3 seconds, which raises the error
rate everywhere, but to set a cut: any item not clearly on its way after 40 seconds is left blank (or guessed, if the
scoring rule makes a guess worth it, \cref{met:iv:online-assessments:timed}), and the time goes to the items that can be
finished. Look at the clock at items 30, 40 and 50 again, not after every item.
\end{example}

\begin{example}[Reading the bounds of a coding item]\label{ex:iv:online-assessments:pairs}
\emph{``Given up to $10^5$ trade prices (integers) and a target, count the pairs $i < j$ whose prices sum to the
target.''} The bound decides the method: $n^2/2 = 5 \times 10^9$ pair checks will not finish in a one-second limit, so
the double loop is only a check. One pass with a hash map of counts of prices seen so far adds, at each price $p$, the
count of $\text{target} - p$ already seen, then records $p$: $O(n)$ time. On the prices 3, 5, 2, 5, 1 with target 7 the
pass adds 0, 0, 1 (the 5 before the 2), 1 (the 2 before the second 5) and 0: two pairs, which the double loop confirms.
The hidden tests to expect: all prices equal (with target twice that price, $n(n-1)/2$ pairs, which overflows 32-bit
integers at $n = 10^5$), no pair, negative prices if the statement allows them, and $n = 1$.
\end{example}

\section{Question bank}

\begin{interviewq}[$\star$ \iqroles{trader} \iqfirm{market maker}]\label{iq:iv:online-assessments:1}
Six seconds each, no calculator: (a) $48 \times 0.35$; (b) $\tfrac78 - \tfrac23$, to four decimals;
(c) $2.4 \div 0.016$.
\end{interviewq}

\begin{interviewq}[$\star$ \iqroles{trader, researcher, developer} \iqfirm{any}]\label{iq:iv:online-assessments:2}
Give the next term: (a) 3, 8, 15, 24, 35; (b) 2, 3, 5, 9, 17; (c) 1, 4, 2, 8, 3, 12, 4.
\end{interviewq}

\begin{interviewq}[$\star$ \iqroles{trader} \iqfirm{proprietary firm}]\label{iq:iv:online-assessments:3}
A test gives $+1$ for a right answer, $-\tfrac12$ for a wrong one and 0 for a blank, on four-option items.
Should you guess blind? After eliminating one option? Two?
\end{interviewq}

\begin{interviewq}[$\star\star$ \iqroles{developer} \iqfirm{market maker}]\label{iq:iv:online-assessments:4}
An automated assessment asks for the length of the longest run of strictly increasing consecutive prices in a
list of up to $10^6$ trade prices. Write it, and list the hidden test cases you expect it to face.
\end{interviewq}

\begin{interviewq}[$\star\star$ \iqroles{trader} \iqfirm{market maker}]\label{iq:iv:online-assessments:5}
Eighty items in eight minutes, $-1$ for a wrong answer. The first sixty are easy: you take 5 seconds each at
97\% accuracy. The last twenty are hard: 12 seconds each at 70\%. Items are answered in order and cannot be
revisited. What is your plan and your expected score? What changes if your accuracy on the hard items is only
45\%?
\end{interviewq}

\begin{interviewq}[$\star\star$ \iqroles{researcher, trader} \iqfirm{any}]\label{iq:iv:online-assessments:6}
Fifteen seconds each: (a) two fair dice; the chance their sum is prime? (b) three fair coins; the chance of
at least two heads?
\end{interviewq}

\begin{interviewq}[$\star\star$ \iqroles{trader, researcher, developer} \iqfirm{any}]\label{iq:iv:online-assessments:7}
Five traders A, B, C, D, E sit in a row of five seats. C sits at the left end. A sits at neither end and not
next to B. D sits immediately to the left of E. B sits somewhere to the right of A. Who sits in the middle?
\end{interviewq}

\begin{interviewq}[$\star\star\star$ \iqroles{trader, risk} \iqfirm{market maker}]\label{iq:iv:online-assessments:8}
A game shows a balloon. Each pump adds one point to a pot; the balloon bursts at a point drawn uniformly from 1
to 64 pumps, and then the pot is lost. You choose in advance how many pumps to make, then bank. How many
maximise the expected pot, and what is it? What might the game be scoring besides the pot?
\end{interviewq}

\begin{interviewq}[$\star\star\star$ \iqroles{developer} \iqfirm{proprietary firm}]\label{iq:iv:online-assessments:9}
Given up to $2 \times 10^5$ trade prices and up to $2 \times 10^5$ queries, each asking how many trades had a
price in $[a, b]$, answer all queries within a one-second limit. What is your approach and its complexity, and
which edge cases must it handle?
\end{interviewq}

\begin{interviewq}[$\star\star\star$ \iqroles{researcher, trader} \iqfirm{any}]\label{iq:iv:online-assessments:10}
A test's score is your ability plus independent noise with standard deviation 5; the pass mark is 60 and your
ability is 57. What is your chance of passing one sitting? Two independent sittings? If practice raises your
ability by a quarter of the spread of candidates' abilities, which is 8, what is your chance at one sitting
then? What does this suggest about when to take a test that cannot be retaken for months?
\end{interviewq}

\begin{omsources}
\item J.~P.~Hausknecht, J.~A.~Halpert, N.~T.~Di Paolo and M.~O.~Moriarty Gerrard, ``Retesting in selection: a
meta-analysis of coaching and practice effects for tests of cognitive ability'', \emph{Journal of Applied
Psychology} 92(2), 2007, 373--385.
\item Equality Act 2010, section 20 and Schedule 8; 29 CFR 1630.11.
\item T.~H.~Cormen, C.~E.~Leiserson, R.~L.~Rivest and C.~Stein, \emph{Introduction to Algorithms}, 4th
edition, MIT Press, 2022, for the complexity rules of thumb.
\end{omsources}
